% answers_short.tex —— 2014 年·山东卷·理科 · 简短答案速查
% 由 tools/make_short_answers.py 自 answers.tex 的 \gksolhead/\gkanswer 提取，
% 并与 2014/2014_数学真题_山东理_答案要点.md 逐题比对后生成，勿手改（改源 MD 后重跑）。

\begin{gknotice}
\textbf{使用说明}：本节只给最终答案便于快速核对，完整推导见解析版。答案为编写者独立推导，\textbf{非官方标准答案}。
\end{gknotice}

\gkpart{选择题}
\begin{center}\begin{tabular}{ccccccccccc}
\toprule
题号 & 1 & 2 & 3 & 4 & 5 & 6 & 7 & 8 & 9 & 10 \\
\midrule
答案 & D & C & C & A & D & D & C & B & B & A \\
\bottomrule
\end{tabular}\end{center}

\gkpart{填空题}
\begin{itemize}[leftmargin=2.6em]
  \item[\textbf{11.}] $3$。$x=1,2,3$ 时条件成立，依次得 $n=1,2,3$；$x=4$ 时 $x^2-4x+3=3>0$，输出 $n=3$。
  \item[\textbf{12.}] $\dfrac16$。由 $|\overrightarrow{AB}||\overrightarrow{AC}|\cos A=\tan A$ 得 $|\overrightarrow{AB}||\overrightarrow{AC}|=\dfrac23$，故 $S=\dfrac12\cdot\dfrac23\cdot\sin\dfrac\pi6=\dfrac16$。
  \item[\textbf{13.}] $\dfrac14$。以 $P$ 为顶点用体积公式，$V_1=\dfrac14V_2$。
  \item[\textbf{14.}] $2$。$x^3$ 项为 $\mathrm C_6^3a^3b^3x^3$，由 $20a^3b^3=20$ 得 $ab=1$，故 $a^2+b^2\geqslant2|ab|=2$。
  \item[\textbf{15.}] $(2\sqrt{10},+\infty)$。由对称定义得 $h(x)=6x+2b-\sqrt{4-x^2}$，条件化为 $b>\sqrt{4-x^2}-3x$ 恒成立，右端最大值为 $2\sqrt{10}$。
\end{itemize}

\gkpart{解答题}
\begin{itemize}[leftmargin=2.6em]
  \item[\textbf{16.}] (1) $m=\sqrt3$，$n=1$。代入两点坐标得关于 $m$，$n$ 的方程组，解得 $m=\sqrt3$，$n=1$；(2) $\left[k\pi-\dfrac\pi2,k\pi\right]$（$k\in\mathbf Z$）。由最高点到 $(0,3)$ 的距离最小值为 $1$ 得 $\varphi=\dfrac\pi6$，$g(x)=2\cos2x$
  \item[\textbf{17.}] (1) 见解析版；(2) $\dfrac{\sqrt5}5$。建系得平面 $C_1D_1M$ 的法向量与底面的法向量，由夹角余弦得 $\dfrac1{\sqrt5}$
  \item[\textbf{18.}] (1) $\dfrac3{10}$。两次落点在乙上的概率分别为 $\dfrac56$，$\dfrac45$，故 $P=\dfrac56\cdot\dfrac15+\dfrac16\cdot\dfrac45=\dfrac3{10}$；(2) $P(\xi=0)=\dfrac1{30}$，$P(\xi=1)=\dfrac16$，$P(\xi=2)=\dfrac15$，$P(\xi=3)=\dfrac2{15}$，$P(\xi=4)=\dfrac{11}{30}$，$P(\xi=6)=\dfrac1{10}$；$E\xi=\dfrac{91}{30}$。逐对枚举两次得分求和，期望可用两球得分期望相加复核
  \item[\textbf{19.}] (1) $a_n=2n-1$。由 $S_2^2=S_1S_4$ 得 $a_1=1$；(2) $T_n=1+\dfrac{(-1)^{n-1}}{2n+1}$。拆项 $b_n=(-1)^{n-1}\left(\dfrac1{2n-1}+\dfrac1{2n+1}\right)$ 后相邻相消
  \item[\textbf{20.}] (1) 单调递减区间为 $(0,2)$，单调递增区间为 $(2,+\infty)$。$k\leqslant0$ 时 $\mathrm e^x-kx>0$，$f'(x)$ 与 $x-2$ 同号；(2) $\left(\mathrm e,\dfrac{\mathrm e^2}2\right)$。$f'(x)=0\iff k=\dfrac{\mathrm e^x}x$，该函数在 $(0,2)$ 上先减后增，最小值 $\mathrm e$、右端点值 $\dfrac{\mathrm e^2}2$
  \item[\textbf{21.}] (1) $y^2=4x$。由正三角形条件 $|AD|=|DF|$ 得 $p^2+4p-12=0$，故 $p=2$；(2) ① 定点为 $(1,0)$；② 存在，最小值 $16$。设 $A(t^2,2t)$，则 $S=2\left|t+\dfrac1t\right|^3\geqslant16$
\end{itemize}
